\chapter{Studiu bibliografic/Stadiul actual în domeniu}

Documentarea bibliografică are ca obiectiv prezentarea stadiului actual al domeniului/ sub-domeniului în care se situează tema, prezentarea cercetărilor similare şi raportarea abordării din lucrare la acestea. 

Acest capitol reprezintă cca. 10--15\% din lucrare.

Referințele se scriu în secțiunea Bibliografie. Formatul referințelor trebuie să fie de tipul IEEE sau asemănător. 

Referințele bibliografice se vor face pentru fiecare carte, articol sau material folosit pentru elaborarea lucrării de disertație. 

În secțiunea Bibliografie sunt exemple de referințe pentru articol la conferințe sau seminarii  \cite{Zhou04}, articol în jurnal \cite{Antoniou04}, sau cărți \cite{strunk}. 

Referințele spre aplicații sau resurse online (pagini de internet) trebuie să includă cel puțin o denumire sugestivă pe lângă link-ul propriu zis \cite{webpage}, plus alte informații dacă sunt disponibile (autori, an, etc.). Referințele care prezintă doar link spre resursa online se vor plasa în footer-ul paginii unde sunt referite.

Citarea referințelor în text este obligatorie, vezi exemplul de mai jos (în funcție de tema proiectului se poate varia modul de prezentare a metodei/aplicației).

\section{Abordări similare}

Citarea referințelor se face ca în exemplele din \ref{subsec:s10}, de mai sus și din citările următoare.

În articolul \cite{Antoniou04} autorul descrie configurația tehnică a unei "honeynet" și prezintă câteva atacuri de actualitate asupra honeynet, precum și o serie de recomandări pentru securizarea sistemelor conectate la rețele de calculatoare.

În capitolul 4 al \cite{strunk}, \dots.







